\documentclass[11pt,a4paper]{coverletter}

\senderblock
  {Radheem Bin Razi}
  {Distributed Systems \& Cloud-Native Engineer}
  {Ilmenau, Germany \textbar\ \href{mailto:sheikh.radheem@gmail.com}{sheikh.radheem@gmail.com} \textbar\ \href{https://radheem.github.io/my-cv}{portfolio}}

\begin{document}

%% ===== ENGLISH =====
\recipient{Randstad Deutschland}{Hiring Team}{}

\opening{Dear Hiring Team,}

Building data pipelines that survive worker crashes and still pass a data-quality gate is the part of engineering I keep returning to, so a Data Engineer role centered on robust ETL/ELT, Python, and Terraform on AWS speaks directly to how I already work. I would welcome the chance to bring that to your team's data platform.

Most of my recent work has been pipeline engineering. In my Second Brain project I built a durable, retryable Hatchet ingestion pipeline — extract, chunk, embed, summarize, index — where every step retries independently and survives worker crashes, with NATS JetStream progress streaming and pgvector semantic search on top. On the IRS platform I built best-effort ETL that routes records into PostgreSQL, Dgraph, and DocumentDB via declarative routing, with a Ginkgo end-to-end data-quality gate guarding correctness before anything ships. Both run cloud-native on Kubernetes, provisioned as code.

At Bluefin Exchange I was the design authority for core services of a high-throughput crypto exchange, where market and account data moved through Kafka into PostgreSQL and DynamoDB. I drove system-design improvements, owned design and code reviews, and ran Terraform-provisioned, Dockerized infrastructure under strict quality and security standards — the same Infrastructure-as-Code discipline this role asks for. Earlier, at Seed Labs, I cleaned, analyzed, and visualized data for stakeholders in Python, translating business questions into concrete data workflows — the requirements-to-solution translation that data engineering lives on.

Python, SQL, ETL design, Terraform, and AWS are core tools for me, and I am comfortable owning scalable, secure data solutions end to end. I am based in Ilmenau, available full-time now and open to working-student or part-time arrangements, and can relocate within Germany in about two to three weeks. I am Blue Card-ready, so a contract is all that is needed — no sponsorship overhead. My working language is English and my German is at A2 and improving. I would be glad to discuss how I can contribute.

\closing{Sincerely,}{Radheem Bin Razi}

\clearpage
\selectlanguage{ngerman}

%% ===== DEUTSCH =====
\recipient{Randstad Deutschland}{Personalabteilung}{}

\opening{Sehr geehrtes Hiring-Team,}

Datenpipelines zu bauen, die Worker-Abstürze überstehen und trotzdem ein Datenqualitäts-Gate bestehen, ist der Teil der Entwicklung, zu dem ich immer wieder zurückkehre. Daher spricht mich eine Rolle als Data Engineer mit Fokus auf robustes ETL/ELT, Python und Terraform auf AWS unmittelbar an, weil sie genau meiner Arbeitsweise entspricht. Ich würde mich freuen, dies in die Datenplattform Ihres Teams einzubringen.

Der Großteil meiner jüngsten Arbeit war Pipeline-Engineering. In meinem Projekt Second Brain habe ich eine robuste, wiederholbare Hatchet-Ingestion-Pipeline gebaut — Extrahieren, Chunking, Embedding, Zusammenfassen, Indexieren — bei der jeder Schritt unabhängig wiederholt wird und Worker-Abstürze übersteht, mit NATS-JetStream-Fortschritts-Streaming und semantischer pgvector-Suche obendrauf. Auf der IRS-Plattform habe ich Best-Effort-ETL gebaut, das Datensätze über deklaratives Routing in PostgreSQL, Dgraph und DocumentDB verteilt, mit einem Ginkgo-End-to-End-Datenqualitäts-Gate, das die Korrektheit absichert, bevor etwas ausgeliefert wird. Beide laufen cloud-nativ auf Kubernetes, bereitgestellt als Code.

Bei Bluefin Exchange war ich die Design-Verantwortung für Kerndienste einer Krypto-Börse mit hohem Durchsatz, in der sich Markt- und Kontodaten über Kafka in PostgreSQL und DynamoDB bewegten. Ich trieb Verbesserungen im Systemdesign voran, verantwortete Design- und Code-Reviews und betrieb Terraform-bereitgestellte, dockerisierte Infrastruktur unter strengen Qualitäts- und Sicherheitsstandards — dieselbe Infrastructure-as-Code-Disziplin, die diese Rolle verlangt. Zuvor habe ich bei Seed Labs Daten für Stakeholder in Python bereinigt, analysiert und visualisiert und geschäftliche Fragestellungen in konkrete Daten-Workflows übersetzt — jene Übersetzung von Anforderung zu Lösung, von der Data Engineering lebt.

Python, SQL, ETL-Design, Terraform und AWS sind für mich zentrale Werkzeuge, und ich übernehme gerne durchgängig die Verantwortung für skalierbare, sichere Datenlösungen. Ich bin in Ilmenau ansässig, ab sofort in Vollzeit verfügbar und offen für Werkstudenten- oder Teilzeitmodelle und kann innerhalb von etwa zwei bis drei Wochen innerhalb Deutschlands umziehen. Ich bin Blue-Card-bereit, sodass lediglich ein Vertrag erforderlich ist — kein Aufwand für ein Sponsoring. Meine Arbeitssprache ist Englisch und mein Deutsch befindet sich auf A2-Niveau und verbessert sich stetig. Ich würde mich freuen, zu besprechen, wie ich beitragen kann.

\closing{Mit freundlichen Grüßen,}{Radheem Bin Razi}

\end{document}
